% Options for packages loaded elsewhere
\PassOptionsToPackage{unicode}{hyperref}
\PassOptionsToPackage{hyphens}{url}
\PassOptionsToPackage{dvipsnames,svgnames,x11names}{xcolor}
\documentclass[
  11pt,
  a4paper,
]{article}
\usepackage{xcolor}
\usepackage[margin=2.6cm]{geometry}
\usepackage{amsmath,amssymb}
\setcounter{secnumdepth}{5}
\usepackage{iftex}
\ifPDFTeX
  \usepackage[T1]{fontenc}
  \usepackage[utf8]{inputenc}
  \usepackage{textcomp} % provide euro and other symbols
\else % if luatex or xetex
  \usepackage{unicode-math} % this also loads fontspec
  \defaultfontfeatures{Scale=MatchLowercase}
  \defaultfontfeatures[\rmfamily]{Ligatures=TeX,Scale=1}
\fi
\usepackage{lmodern}
\ifPDFTeX\else
  % xetex/luatex font selection
  \setmainfont[]{Libertinus Serif}
  \setmonofont[Scale=0.85]{Latin Modern Mono}
  \setmathfont[]{Latin Modern Math}
\fi
% Use upquote if available, for straight quotes in verbatim environments
\IfFileExists{upquote.sty}{\usepackage{upquote}}{}
\IfFileExists{microtype.sty}{% use microtype if available
  \usepackage[]{microtype}
  \UseMicrotypeSet[protrusion]{basicmath} % disable protrusion for tt fonts
}{}
\makeatletter
\@ifundefined{KOMAClassName}{% if non-KOMA class
  \IfFileExists{parskip.sty}{%
    \usepackage{parskip}
  }{% else
    \setlength{\parindent}{0pt}
    \setlength{\parskip}{6pt plus 2pt minus 1pt}}
}{% if KOMA class
  \KOMAoptions{parskip=half}}
\makeatother
\usepackage{longtable,booktabs,array}
\newcounter{none} % for unnumbered tables
\usepackage{calc} % for calculating minipage widths
% Correct order of tables after \paragraph or \subparagraph
\usepackage{etoolbox}
\makeatletter
\patchcmd\longtable{\par}{\if@noskipsec\mbox{}\fi\par}{}{}
\makeatother
% Allow footnotes in longtable head/foot
\IfFileExists{footnotehyper.sty}{\usepackage{footnotehyper}}{\usepackage{footnote}}
\makesavenoteenv{longtable}
\setlength{\emergencystretch}{3em} % prevent overfull lines
\providecommand{\tightlist}{%
  \setlength{\itemsep}{0pt}\setlength{\parskip}{0pt}}

\providecommand{\E}{\mathbb{E}}
\providecommand{\N}{\mathbb{N}}
\providecommand{\Z}{\mathbb{Z}}
\providecommand{\R}{\mathbb{R}}
\providecommand{\eps}{\varepsilon}
\providecommand{\ceil}[1]{\left\lceil #1\right\rceil}
\providecommand{\floor}[1]{\left\lfloor #1\right\rfloor}
\AtBeginDocument{\let\setminus\smallsetminus}
\usepackage[most]{tcolorbox}
\newtcolorbox{warning}{colback=red!5,colframe=red!65!black,boxrule=0.8pt,arc=2pt,left=6pt,right=6pt,top=5pt,bottom=5pt}
\usepackage{etoolbox}
\AtBeginEnvironment{longtable}{\small}
\setlength{\emergencystretch}{3em}
\usepackage{fancyhdr}
\pagestyle{fancy}
\fancyhf{}
\fancyhead[L]{\small\itshape Candidate result geodesic\_cycles\_and\_tuttes\_theorem --- GPT-6 Astra, awaiting review}
\fancyhead[R]{\small\thepage}
\renewcommand{\headrulewidth}{0.2pt}
\usepackage{bookmark}
\IfFileExists{xurl.sty}{\usepackage{xurl}}{} % add URL line breaks if available
\urlstyle{same}
\hypersetup{
  pdftitle={A candidate counterexample to Geodesic cycles and Tutte's Theorem},
  pdfauthor={github.com/graph-theory-AI --- machine-generated candidate writeup},
  colorlinks=true,
  linkcolor={blue!50!black},
  filecolor={Maroon},
  citecolor={Blue},
  urlcolor={blue!50!black},
  pdfcreator={LaTeX via pandoc}}

\title{A candidate counterexample to Geodesic cycles and Tutte's
Theorem}
\author{\href{https://github.com/graph-theory-AI}{github.com/graph-theory-AI}
--- machine-generated candidate writeup}
\date{16 September 2026}

\begin{document}
\maketitle
\begin{abstract}
The eight-vertex graph obtained by inserting a degree-three vertex into
every face of a tetrahedron is a counterexample. This candidate disproof
was generated by GPT-6 Astra in a single model pass for the Graph Theory
AI project. It was selected from the campaign because the model marked
it \texttt{would\_publish:\ true}; it has not been checked by the
adversarial referee, a human mathematician, or a novelty search.
\end{abstract}

\section{The problem and the claimed
result}\label{the-problem-and-the-claimed-result}

This document concerns
\href{https://graph-theory-ai.github.io/graph-conjectures/op/geodesic_cycles_and_tuttes_theorem/}{Geodesic
cycles and Tutte's Theorem} from
\emph{http://www.openproblemgarden.org/op/geodesic\_cycles\_and\_tuttes\_theorem}.

\subsection{Problem statement}\label{problem-statement}

Problem If \(G\) is a \(3\) -connected finite graph, is there an
assignment of lengths \(\ell: E(G) \to \mathbb{R}^+\) to the edges of
\(G\) , such that every \(\ell\) -geodesic cycle is peripheral ?

\subsection{Catalog context}\label{catalog-context}

A cycle \(C\) is \(\ell\) -geodesic if for every two vertices \(x,y\) on
\(C\) there is no \(x\) - \(y\) \textasciitilde path in \(G\) shorter,
with respect to \(\ell\) , than both \(x\) - \(y\) \textasciitilde arcs
on \(C\) . It is not hard to prove {[}GS{]} that for every finite graph
\(G\) and every assignment of edge lengths
\(\ell: E(G) \to \mathbb{R}^+\) the \(\ell\) -geodesic cycles of \(G\)
generate its cycle space. Thus, a positive answer to the problem would
imply a new proof of Tutte's classical theorem {[}T{]} that the
peripheral cycles of a \(3\) -connected finite graph generate its cycle
space.

\subsection{Claimed result}\label{claimed-result}

\textbf{Model verdict: disproved} (confidence: high).

The eight-vertex graph obtained by inserting a degree-three vertex into
every face of a tetrahedron is a counterexample.

\emph{The model's own caveat:} Peripheral has its standard induced,
vertex-nonseparating meaning; bibliographical novelty has not been
checked.

\section{Model writeup}\label{model-writeup}

The text below is the model output verbatim, apart from moving its
machine-readable verdict block into the summary above and shifting
Markdown heading levels for this document. Mechanical TeX defects and
equation tags were normalized where needed for compilation.

\subsection{An eight-vertex
counterexample}\label{an-eight-vertex-counterexample}

A cycle is \textbf{peripheral} if it is induced and deleting its
vertices leaves a connected graph.

Let \[
B=\{b_1,b_2,b_3,b_4\},\qquad S=\{s_1,s_2,s_3,s_4\},
\] and define \(G\) by \[
V(G)=B\cup S,\qquad
E(G)=\{b_ib_j:i<j\}\ \cup\ \{s_i b_j:i\ne j\}.
\] Thus \(B\) induces a \(K_4\), the set \(S\) is independent, and
\(s_i\) is adjacent to the three vertices of \(B\setminus\{b_i\}\).

Equivalently, start with a tetrahedron embedded on the sphere and insert
one vertex into each face, joining it to that face's three vertices.

\textbf{Theorem.} For every assignment \(\ell:E(G)\to\mathbb R_{>0}\),
some \(\ell\)-geodesic cycle of \(G\) is not peripheral.

\subsubsection{\texorpdfstring{1. Connectivity and peripheral cycles of
\(G\)}{1. Connectivity and peripheral cycles of G}}\label{connectivity-and-peripheral-cycles-of-g}

The graph \(G\) is 3-connected. Indeed, after deleting at most two
vertices, at least two vertices of \(B\) remain and induce a connected
clique. Every surviving \(s_i\) retains at least one of its three
neighbours in \(B\), so the remaining graph is connected.

There are eight vertices and eighteen edges, hence \[
\dim_{\mathbb F_2}\mathcal C(G)=18-8+1=11.
\]

We next classify the peripheral cycles.

Every induced cycle of \(G\) is a triangle. To see this, an induced
cycle of length at least four cannot contain \(s_i\): its two neighbours
along the cycle would belong to \(B\), and the edge between them would
be a chord. A cycle of length at least four entirely within \(B\) also
has a chord.

The triangles are of two types:

\begin{itemize}
\item
  The four triangles \[
  T_i=G[B\setminus\{b_i\}]
  \] are \textbf{not peripheral}. In \(G-V(T_i)\), the vertex \(s_i\) is
  isolated, whereas the other remaining vertices form a nonempty
  connected component.
\item
  The twelve triangles \[
  P_{i;jk}=s_i b_j b_k s_i,
  \qquad j<k,\quad i\notin\{j,k\},
  \] \textbf{are peripheral}. After deleting such a triangle, two
  adjacent vertices of \(B\) remain. Every other surviving vertex of
  \(S\) has a neighbour among these two vertices.
\end{itemize}

Consequently, \(G\) has exactly twelve peripheral cycles, all triangles.
Each contains \textbf{exactly one edge of the core \(K_4\)}.

\subsubsection{2. Two metric facts}\label{two-metric-facts}

We use two elementary facts, including their proofs to handle arbitrary
positive lengths and possible ties.

\paragraph{Fact 1: Geodesic cycles generate the cycle
space}\label{fact-1-geodesic-cycles-generate-the-cycle-space}

Let \(H\) be any finite graph with positive edge lengths. If a cycle
\(C\) is not geodesic, there is a path \(P\) between vertices
\(x,y\in V(C)\) with \[
\ell(P)<d_C(x,y),
\] where \(d_C\) is distance measured within \(C\).

There is then a path \(Q\), with distinct endpoints \(u,v\in V(C)\) and
all internal vertices outside \(C\), that is shorter than both
\(u\)-\(v\) arcs of \(C\). Indeed, otherwise every excursion of \(P\)
outside \(C\) could be replaced by a no-longer arc of \(C\). This would
produce an \(x\)-\(y\) walk in \(C\) of length at most
\(\ell(P)<d_C(x,y)\), a contradiction.

Write the two \(u\)-\(v\) arcs as \(A,B\). The cycles \[
C_1=A\cup Q,\qquad C_2=B\cup Q
\] satisfy \[
E(C)=E(C_1)\mathbin{\triangle}E(C_2),
\qquad
\ell(C_1)<\ell(C),\quad \ell(C_2)<\ell(C).
\] Induction through the finitely many cycle lengths proves that the
geodesic cycles generate \(\mathcal C(H)\).

\paragraph{Fact 2: An admissible assignment can be perturbed to have
unique shortest
paths}\label{fact-2-an-admissible-assignment-can-be-perturbed-to-have-unique-shortest-paths}

Suppose an assignment has the property that every geodesic cycle is
peripheral. For each nonperipheral cycle, choose one path witnessing its
failure to be geodesic. Each witness gives two strict linear
inequalities: the path is shorter than each of the two corresponding
cycle arcs.

There are finitely many nonperipheral cycles. Thus these inequalities,
and positivity of all edge lengths, persist throughout an open
neighbourhood of the assignment.

Inside this neighbourhood, choose lengths avoiding all equalities \[
\ell(P)=\ell(Q)
\] between distinct simple paths with the same endpoints. These
equalities form finitely many proper hyperplanes. The perturbed
assignment still has no nonperipheral geodesic cycle, and every pair of
vertices now has a unique shortest path.

Importantly, this argument preserves the strict shortcuts for
nonperipheral cycles; it does \textbf{not} assume that particular
peripheral cycles remain geodesic.

\subsubsection{3. Contradiction}\label{contradiction}

Suppose that \(G\) admits a positive length assignment for which every
geodesic cycle is peripheral. By Fact 2, assume additionally that
shortest paths between vertices are unique.

By Fact 1, the geodesic cycles generate the 11-dimensional cycle space.
Since all of them belong to the twelve peripheral triangles listed
above, at least eleven of those triangles must be geodesic. Therefore \[
\boxed{\text{At most one peripheral triangle is not geodesic.}}
\qquad\text{(1)}
\]

Call an edge \(uv\) \textbf{tight} if \[
\ell(uv)=d_\ell(u,v).
\] Every edge satisfies \(\ell(uv)\ge d_\ell(u,v)\), so a non-tight edge
has strictly greater length than the distance between its endpoints.

Consider a core edge \(e=b_jb_k\). Let
\(\{p,q\}=\{1,2,3,4\}\setminus\{j,k\}\). Exactly two peripheral
triangles contain \(e\): \[
P_{p;jk}\quad\text{and}\quad P_{q;jk}.
\]

If \(e\) is non-tight, these two triangles cannot both be geodesic.
Otherwise, geodesicity for the pair \(b_j,b_k\) would imply \[
\ell(b_js_p)+\ell(s_pb_k)
=d_\ell(b_j,b_k)
=\ell(b_js_q)+\ell(s_qb_k).
\] The two distinct paths \(b_js_pb_k\) and \(b_js_qb_k\) would both be
shortest, contradicting uniqueness.

Thus every non-tight core edge belongs to a non-geodesic peripheral
triangle. Different core edges require different such triangles, because
each peripheral triangle contains exactly one core edge. By (1), \[
\boxed{\text{At most one core edge is non-tight.}}
\qquad\text{(2)}
\]

There is consequently a core triangle \(T_i\) containing no non-tight
edge: if the exceptional edge exists, choose a core triangle avoiding
it.

All three edges of this \(T_i\) are tight. For every pair of vertices of
\(T_i\), their joining edge therefore realizes their distance in \(G\).
Hence \(T_i\) is geodesic.

But \(T_i\) is not peripheral, as established in Section 1. This is the
required contradiction. \(\square\)

\subsubsection{Completeness and status}\label{completeness-and-status}

The perturbation argument covers assignments with tied shortest paths as
well as generic assignments. Thus the counterexample rules out
\textbf{every} positive edge-length assignment, not merely those in
general position.

The graph is simple, planar, and 3-connected, so the assertion fails
even within the planar subclass. The proof is self-contained; I have not
checked whether this counterexample or argument has previously appeared
in the literature.

\appendix

\section{Provenance and review
status}\label{provenance-and-review-status}

\begin{warning}

\textbf{UNREFEREED MODEL OUTPUT.} This document was selected solely
because GPT-6 Astra labelled its own result
\texttt{would\_publish:\ true} and returned \texttt{proved} or
\texttt{disproved}. It has not passed the adversarial LLM referee used
for the earlier Sol campaign, has not been checked by a human
mathematician, and has not been checked for novelty. Treat every
mathematical and bibliographic claim as unverified.

\end{warning}

{\def\LTcaptype{none} % do not increment counter
\begin{longtable}[]{@{}
  >{\raggedright\arraybackslash}p{(\linewidth - 2\tabcolsep) * \real{0.5000}}
  >{\raggedright\arraybackslash}p{(\linewidth - 2\tabcolsep) * \real{0.5000}}@{}}
\toprule\noalign{}
\endhead
\bottomrule\noalign{}
\endlastfoot
Catalog id & \texttt{geodesic\_cycles\_and\_tuttes\_theorem} \\
Catalog entry &
\href{https://graph-theory-ai.github.io/graph-conjectures/op/geodesic_cycles_and_tuttes_theorem/}{Geodesic
cycles and Tutte's Theorem} \\
Source corpus & opg \\
Campaign leg & OpenProblemGarden first pass (\texttt{attacks\_opg}) \\
Source paper / entry &
http://www.openproblemgarden.org/op/geodesic\_cycles\_and\_tuttes\_theorem \\
Model verdict & \textbf{disproved} (confidence: high) \\
Selection criterion & \texttt{would\_publish:\ true} and verdict
\texttt{proved} or \texttt{disproved} \\
Model & \texttt{gpt-6-astra}, reasoning effort \texttt{max},
\texttt{mode=pro}, flex service tier \\
Original artifact &
\texttt{attacks\_opg/geodesic\_cycles\_and\_tuttes\_theorem/output.md} \\
Independent review & \textbf{None} \\
Model's caveats & Peripheral has its standard induced,
vertex-nonseparating meaning; bibliographical novelty has not been
checked. \\
\end{longtable}
}

The generated Markdown and LaTeX sources for this PDF are stored under
\texttt{to\_review\_astra/src/geodesic\_cycles\_and\_tuttes\_theorem/}.

\end{document}
